% ============================================================================
%  PHẦN ĐẦU: Lời cảm ơn · Tóm tắt · Mục lục · Danh sách hình/bảng · Từ viết tắt
% ============================================================================

\section*{Lời cảm ơn}
\addcontentsline{toc}{section}{Lời cảm ơn}

Nhóm 7 xin gửi lời cảm ơn chân thành đến PGS.TS. Nguyễn Đình Hiển, giảng viên môn
\emph{Trí tuệ nhân tạo} (CS106), đã giới thiệu các phương pháp biểu diễn tri thức và suy
diễn làm nền tảng cho đồ án, cũng như dòng nghiên cứu mô hình tri thức cho văn bản pháp
luật mà hệ thống của nhóm kế thừa.

Nhóm cũng cảm ơn cộng đồng mã nguồn mở Python, scikit-learn, FastAPI và Pydantic đã cung
cấp hạ tầng để một nhóm sinh viên có thể xây dựng, kiểm thử và đóng gói trọn vẹn một hệ
tra cứu pháp luật chạy được trên máy tính cá nhân.

\vspace{1em}
\hfill Thành phố Hồ Chí Minh, tháng 9 năm 2026

\hfill Nhóm 7

\clearpage
% ----------------------------------------------------------------------------
\section*{Tóm tắt}
\addcontentsline{toc}{section}{Tóm tắt}

Đồ án xây dựng một \textbf{hệ thống tra cứu kiến thức pháp luật giao thông đường bộ}
theo khung pháp lý 2024--2026, gồm Luật Trật tự, an toàn giao thông đường bộ số
36/2024/QH15, Nghị định 168/2024/NĐ-CP và hai văn bản sửa đổi là Luật 118/2025/QH15 và
Nghị định 238/2026/NĐ-CP. Khác với hướng truy hồi văn bản hay sinh câu trả lời bằng mô
hình ngôn ngữ lớn, hệ thống biểu diễn tri thức \emph{tường minh} bằng mô hình
$K = (C, R, \mathit{Rules}, F, \mathit{Keyphrase})$ và \emph{không sinh ngôn ngữ tự
nhiên}: mọi câu trả lời đều là nguyên văn điều khoản kèm căn cứ tới mức điều -- khoản --
điểm.

Cơ sở tri thức gồm 73 khái niệm, 482 quan hệ, 109 quy tắc, 345 hành vi vi phạm và 1.674
cụm từ khoá; 100\% mục tri thức có căn cứ pháp lý và khoảng hiệu lực, 63 sửa đổi được lưu
như dữ liệu để trả lời được câu hỏi ``tại thời điểm $t$ thì áp dụng quy định nào''. Thuật
giải xử lý truy vấn sáu bước (B1--B6) phân loại câu hỏi vào bảy lớp bài toán, xếp hạng tri
thức bằng hàm điểm lai ba thành phần (keyphrase, ngữ nghĩa TF-IDF, ngữ cảnh) và dùng bộ
suy diễn số học để chọn đúng khung phạt theo ngưỡng nồng độ cồn, tốc độ.

Trên bộ đánh giá 120 câu hỏi có đáp án chuẩn gán tới định danh tri thức, hệ thống đạt độ
chính xác phân lớp \textbf{95,83\%}, \textbf{Top-1 76,67\%} (khoảng tin cậy 95\%:
68,3--83,3\%), Top-5 95,00\% và MRR 0,8384, với thời gian trả lời trung bình 6,2~ms. Thí
nghiệm loại bỏ thành phần cho thấy suy diễn số học nâng Top-1 của nhóm câu có giá trị số
từ 40,9\% lên 95,5\%, và kiểm định McNemar ($p = 0{,}000488$) khẳng định cải thiện này có ý
nghĩa thống kê. Mọi hạn chế --- trong đó có việc cấu hình mặc định chưa từ chối được truy
vấn ngoài lĩnh vực --- đều được đo và khoá lại bằng kiểm thử. Toàn bộ số liệu tái lập được
bằng một lệnh chạy.

\vspace{0.6em}
\noindent\textbf{Từ khoá:} biểu diễn tri thức, hệ tra cứu pháp luật, giao thông đường bộ,
keyphrase, suy diễn số học, hiệu lực theo thời gian, tiếng Việt.

\clearpage
% ----------------------------------------------------------------------------
{\setstretch{1.05}
\tableofcontents
\clearpage
\listoffigures
\addcontentsline{toc}{section}{Danh sách hình}
\vspace{1.2em}
\listoftables
\addcontentsline{toc}{section}{Danh sách bảng}
}

\clearpage
% ----------------------------------------------------------------------------
\section*{Danh mục từ viết tắt}
\addcontentsline{toc}{section}{Danh mục từ viết tắt}

\begin{center}
\renewcommand{\arraystretch}{1.15}
\begin{tabularx}{\textwidth}{@{}l L@{}}
\toprule
\textbf{Viết tắt} & \textbf{Ý nghĩa} \\
\midrule
GTĐB      & Giao thông đường bộ \\
GPLX      & Giấy phép lái xe \\
NĐ-CP     & Nghị định của Chính phủ \\
QH        & Luật do Quốc hội ban hành \\
VBHN-VPQH & Văn bản hợp nhất của Văn phòng Quốc hội \\
KB        & Knowledge Base --- cơ sở tri thức \\
TF-IDF    & Term Frequency -- Inverse Document Frequency \\
MRR       & Mean Reciprocal Rank --- hạng nghịch đảo trung bình \\
KTC       & Khoảng tin cậy \\
AUC       & Area Under the ROC Curve --- diện tích dưới đường cong ROC \\
RAG       & Retrieval-Augmented Generation --- sinh văn bản có truy hồi \\
LLM       & Large Language Model --- mô hình ngôn ngữ lớn \\
API       & Application Programming Interface \\
CI        & Continuous Integration --- tích hợp liên tục \\
xfail     & Kiểm thử được đánh dấu ``dự kiến thất bại'' kèm số đo hạn chế \\
\bottomrule
\end{tabularx}
\end{center}
